\subsection{Comments on the Fayet-Iliopoulos Term in Field Theory and Supergravity --- arXiv:0904.1159}

\textbf{Authors:} Zohar Komargodski, Nathan Seiberg

\paragraph{主な主張}
場に依存しないFI項について、$\xi\ll M_{\mathrm P}^{2}$の超重力記述では厳密な連続大域対称性が必要となり、その存在を禁じる量子重力の原則と両立しないと論じた\cite{Komargodski:2009pc}。

\paragraph{新規性}
超カレント多重項の超場ゲージ変換に対する非不変性から、FI項の繰り込み制約と超重力への結合の障害を統一的に説明した。

\paragraph{位置付け}
定数FI項を用いたSUSY破れ・宇宙論模型の整合性を検討する基礎文献であり、場に依存する有効D項とは区別して読む。

\paragraph{コメント（読書メモ）}
檜垣さんの説明を記録した読書メモ：弦理論でゲージ化された$U(1)_R$が見えない理由を考えるうえで、この論文が有力な議論になるとのこと。定数FI項がある場合、エネルギー運動量テンソルやsupercurrentは超場の全ゲージ変換に対して不変ではなく、対応する保存電荷が不変でも、カレントを通じた超重力への結合には問題が生じる、と理解した。厳密な大域対称性があれば結合できるが、そのような対称性は量子重力には恐らく存在しない、という説明だった。読書時にはこれを「$U(1)_R$はゲージ化されない」とまとめていたが、本論文の結論は上記の定数FI項と仮定の範囲で読む。場に依存するD項と定数FI項の違いについて、Grimm--Louisの式(4.21)\cite{Grimm:2004uq}との対応を確認したい。SUSY breakingの関連文献。
